\section{The generator}\label{sec:generator}

We can now set up the machine of \S\ref{sec:intro} on a small scale.

\subsection{The database of arches}

Our arches are the five bridges, all proved in Lean in \lp{}:

\begin{center}\small
\begin{tabularx}{\linewidth}{l>{\raggedright\arraybackslash}X>{\raggedright\arraybackslash}X}
\toprule
& invariant of the topos & reading on the shape $C$\\
\midrule
A & Boolean (excluded middle) & $C$ is a groupoid\\
B & De Morgan (weak excluded middle) & $C$ satisfies the Ore condition\\
C & the topos itself, up to equivalence & $\Kar(C)$ up to equivalence\\
D & essential points & objects of $\Kar(C)$\\
E & lattice of subtopoi & topologies finer than $J$\\
\bottomrule
\end{tabularx}
\end{center}

To these we add one arch from topology, the characterisation of De Morgan sheaf topoi on a space by extremal disconnectedness (\S\ref{sec:sites}).

\subsection{The database of decks}

A deck is a topos with two presentations, at least one of which is familiar. Here are ours; each line says \q{the structures on the left are the presheaves on the shape on the right}.

\begin{center}\small
\begin{tabularx}{\linewidth}{>{\raggedright\arraybackslash}X>{\raggedright\arraybackslash}p{5cm}l}
\toprule
structures & shape $C$ (topos $\PSh(C)$) & see\\
\midrule
right $M$-sets; $G$-sets & $\B M$; $\B G$ & Ex.~\ref{ex:Msets}\\
deterministic automata over $\Sigma$ & $\B(\Sigma^*)$ & Ex.~\ref{ex:automata}\\
discrete dynamical systems & $\B(\mathbb N,+)$ & Ex.~\ref{ex:automata}\\
sheaves on the Alexandrov space $X_P$ & $P^{\op}$ & Prop.~\ref{prop:alexandrov}\\
families of sets along a rewriting system & $(S,\to^*)^{\op}$ & \S\ref{sec:preorders}\\
directed graphs; reflexive graphs & $\mathsf G$; $\mathsf G_r$ & Ex.~\ref{ex:graphs}\\
simplicial, augmented, symmetric simplicial sets & $\Delta$, $\Delta_+$, $\mathsf{Fin}_{>0}$ & \S\ref{sec:shapes}\\
the topos of trees & $\omega$ & Ex.~\ref{ex:preorderpsh}\\
$M_n(k)$-sets, $T_n$-sets & $\B M_n(k)$, $\B T_n$ & \S\ref{sec:points}\\
models of $\mathbb T$ in topoi ($\mathbb T$ Horn) & $\fp\mathbb T\text{-mod}^{\op}$ & Ex.~\ref{ex:classifying}\\
$G$-torsors in topoi & $\B G$ & Ex.~\ref{ex:classifying}\\
\bottomrule
\end{tabularx}
\end{center}

Bridge~\ref{br:morita} adds decks of a second kind: two shapes $C$, $D$ with $\Kar(C)\simeq\Kar(D)$, such as a monoid $M$ and its \q{corner} $eMe$ (\S\ref{sec:monoids}).

\subsection{Running the generator}

For each deck and each arch, the generator writes down the reading of the arch on both sides of the deck, and equates them. Mechanically, this produces statements of the form
\[
\text{(familiar property of the structures)}\iff\text{(combinatorial property of the shape)}.
\]
The left side is obtained by \emph{unfolding} the definition of the invariant in the familiar language (Lemma~\ref{lem:neg} does this for negation); the right side is obtained by \emph{computing} the arch on the shape (checking the Ore condition on $\B M$, for instance). The two steps that Caramello's talk singles out as hard to automate are exactly these: choosing which decks are worth trying, and rewriting the unfolded statement \q{into human readable statements} \cite[slide 17]{caramello-mines}. We did both by hand.

\subsection{How we present each output}

In \S\ref{sec:monoids}--\S\ref{sec:theories}, each generated statement is presented in the same way:
\begin{itemize}
\item \textbf{the statement}, in the language of the familiar structures;
\item \textbf{in words}, what it means;
\item \textbf{how the generator finds it}: which deck, which arch;
\item \textbf{a proof without topoi}, which a reader can check without trusting anything above;
\item \textbf{examples};
\item \textbf{is it known?}, what our literature search found (details in \S\ref{sec:lit}).
\end{itemize}
